\documentclass[11pt]{article}
\usepackage[T1]{fontenc}
\usepackage[margin=1in]{geometry}
\usepackage{amsmath,amssymb,amsthm,mathtools,booktabs,graphicx,hyperref}
\setlength{\emergencystretch}{3em}
\newtheorem{theorem}{Theorem}
\newtheorem{lemma}[theorem]{Lemma}
\newtheorem{definition}[theorem]{Definition}
\newcommand{\SAT}{\operatorname{SAT}}
\newcommand{\gammats}{2\cos(\pi/7)}
\title{Main 7: Exact Fingerprints and Adaptive Hashing in the WRAPPER}
\author{PNP research working manuscript}
\date{7 October 2026}
\begin{document}
\maketitle

\begin{abstract}
This manuscript asks whether changing a WRAPPER hash from four to six or
eight displayed digits can supply the missing size reduction.  Exact replay
after a hash match makes correctness independent of the hash length, but it
also makes the number of retained semantic classes independent of the hash
length.  Omitting replay turns collision soundness into the missing semantic
merge theorem.  Shorter fingerprints can improve finite storage constants;
longer or adaptive fingerprints can reduce expected collision work.  Neither
gives the deterministic all-input contraction required by the current
Williams endpoint.  We state the exact ledgers, recover the project's earlier
digest and bitstream experiments, and record a collision-safe adaptive cache
as useful infrastructure.  No proof that $P\ne NP$ is claimed.
\end{abstract}

\section{Question and conventions}
A displayed hexadecimal digit carries four bits.  Thus four, six and eight
hexadecimal digits mean respectively $h=16,24,32$ bits.  The existing Python
prototype instead uses \texttt{digest\_bytes=8}, which means $64$ bits or
sixteen hexadecimal digits.  If ``digit'' means a decimal character, the
corresponding information is $h=d\log_2 10$ bits; this changes the numerical
thresholds but none of the theorems below.

Let $X_1,\ldots,X_Q$ be WRAPPER states generated in one realization.  Let
$\equiv$ be the semantic congruence that is sufficient for every future use
of a state, rather than equality of the current SAT answer alone.  Write $R$
for the number of represented $\equiv$-classes.  An $h$-bit fingerprint is a
map
\[
 H:\mathcal X\longrightarrow\{0,1\}^h,
 \qquad B=2^h.
\]
The current exact cache uses $H$ only to select candidates and then compares
the normalized residual streams.  Consequently a collision can increase
work but cannot change the SAT answer.

\section{Exact hashes do not create semantic contraction}
\begin{theorem}[HASH-INDEX-NO-SHRINK1]
Suppose cache reuse occurs only after an exact verifier proves
$X_i\equiv X_j$.  If the verifier is complete, exact deduplication retains
one representative per semantic class and therefore
\[
 q_{\mathrm{out}}=R.
\]
If the verifier is sound but incomplete, then $q_{\mathrm{out}}\ge R$.
Neither conclusion depends on $h$.
\end{theorem}
\begin{proof}
Soundness prevents representatives of two different congruence classes from
being identified.  Completeness identifies every later representative of a
class with the first one, leaving exactly one per class.  An incomplete test
may leave additional representatives.  The fingerprint changes which pairs
are sent to the verifier; it does not change the equivalence relation proved
by that verifier.
\end{proof}

If equal hashes are merged without verification, correctness requires
\[
 H(x)=H(y)\Longrightarrow x\equiv y.
\tag{1}
\]
Distinct semantic classes must then have distinct hash values, so
\[
 2^h\ge R,\qquad h\ge\lceil\log_2R\rceil.
\tag{2}
\]
Proving (1) for a cheaply constructed small-range map is already the desired
semantic quotient theorem.  A one-bit example exposes the acquisition issue:
$H(F)=\SAT(F)$ perfectly separates the two current-answer classes, but
computing it is SAT itself.  It also does not justify reuse under a richer
future interface.

Input-dependent perfect hashing does not evade the theorem.  It can assign
short collision-free labels to the $Q$ syntactic keys after inspecting them,
but it still indexes $Q$ keys and retains $R$ semantic classes.  Its
construction, representatives and source access must be charged.  The
classical universal and perfect hashing results provide strong dictionary
bounds after the key set is fixed; they do not supply a SAT congruence
\cite{carterwegman1979,fks1984}.

\section{Collision-replay ledger}
Suppose the $Q$ states are pairwise inequivalent and bucket $z$ contains
$m_z$ states.  In a chained exact cache, the number of unequal candidate
comparisons is
\[
 P(H)=\sum_{z=1}^{B}\binom{m_z}{2}.
\tag{3}
\]

\begin{lemma}[Balanced and hostile buckets]
For every $h$-bit assignment of $Q$ states to buckets,
\[
 P(H)\ge
 \max\left\{0,\frac{Q^2/B-Q}{2}\right\}.
\tag{4}
\]
The maximum is $\binom Q2$ and occurs when all states collide.
\end{lemma}
\begin{proof}
Convexity of $x\mapsto\binom{x}{2}$ makes balanced bucket sizes minimize
(3).  Cauchy--Schwarz gives $\sum_zm_z^2\ge Q^2/B$, and substituting into
$P=(\sum_zm_z^2-Q)/2$ gives (4).  Putting every state in one bucket gives the
maximum.
\end{proof}

If one normalized-residual replay costs $E(M)$, collision verification costs
$P(H)E(M)$.  A bounded cache of $K$ representatives changes the hostile
upper bound to $O(QK E(M))$, but a cyclic sequence of $K+1$ distinct states
has no persistent hit under ordinary eviction.  Correctness survives because
misses recompute; no global shrinkage follows.

For a fixed public map on $L$-bit descriptions, one fiber has at least
$2^{L-h}$ members.  Hence a $Q$-state all-collision family exists whenever
\[
 L-h\ge\log_2Q.
\tag{5}
\]
More strongly, after fixing any common prefix of length
$p\le L-h-\log_2Q$, one bucket still contains $Q$ extensions of that prefix.
A streaming exact comparison can then be forced to inspect $\Omega(p)$ bits
for every rejection.  This is a lower bound on this fixed-hash cache design,
not on arbitrary SAT algorithms or on every possible dictionary.

\section{Randomized ledger and the exponent calculation}
For independent uniform $h$-bit fingerprints,
\[
 \mathbb E[P]=\frac{\binom Q2}{2^h},
\qquad
 \Pr[\text{no collision}]
 =\prod_{i=0}^{Q-1}\left(1-\frac{i}{2^h}\right)
 \le \exp\left(-\frac{Q(Q-1)}{2^{h+1}}\right).
\tag{6}
\]
These are the standard pair-collision and birthday calculations behind
universal hashing and rolling fingerprints \cite{carterwegman1979,karprabin1987}.
A union-bound guarantee of collision probability at most $\delta$ needs
approximately
\[
 h\ge2\log_2Q+\log_2(1/\delta)-O(1).
\tag{7}
\]

If
\[
 Q=M^{u+o(1)},\qquad E(M)=M^{v+o(1)},\qquad
 h=\kappa\log_2M,
\]
then reading or hashing every state costs at least $M^{u+v+o(1)}$ in this
streaming model, while expected collision replay has exponent
$2u-\kappa+v$.  Thus
\[
 C_{\mathrm{expected}}
 \ge\max\{u+v,2u-\kappa+v\}.
\tag{8}
\]
At the clean COVER-SHRINK1 target $u=1/3$ and $v=2/3$,
\[
 C_{\mathrm{expected}}\ge\max\{1,4/3-\kappa\}.
\tag{9}
\]
A constant-length hash has $\kappa=0$ and therefore $4/3$ expected replay
exponent in this model.  Choosing $h=(1/3)\log_2M$ reduces expected replay to
the linear scan exponent one.  Choosing $h=(2/3)\log_2M$ reduces the expected
number of false-collision pairs to $O(1)$.  The tag storage
$Qh=M^{1/3+o(1)}$ remains in the same polynomial exponent class.

Equation (9) is an implementation ledger, not the missing contraction.  A
fixed public hash still admits the hostile all-collision family, giving time
$Q^2E=M^{4/3-o(1)}$.  Combined with $M^{2/3}$ workspace, this reaches joint
exponent two and misses the threshold $\gammats\approx1.80194$.  Fresh random
seeds give expected randomized bounds.  The endpoint used in the proof spine
is deterministic, and derandomizing under $P=NP$ does not automatically
preserve these fine-grained exponents.

\section{What four, six and eight hexadecimal digits change}
Table~\ref{tab:digits} uses a uniform-hash model.  The square-root column is
the conventional birthday scale $\sqrt{2^h}$; the probability of a collision
is already a constant near that population.  The existing eight-byte setting
is included to prevent confusing bytes with displayed hexadecimal digits.

\begin{table}[ht]
\centering
\begin{tabular}{rrrrr}
\toprule
display & bits $h$ & buckets $2^h$ & birthday scale & meaning\\
\midrule
4 hex digits & 16 & $65{,}536$ & $256$ & finite toy caches\\
6 hex digits & 24 & $16{,}777{,}216$ & $4{,}096$ & larger finite caches\\
8 hex digits & 32 & $4{,}294{,}967{,}296$ & $65{,}536$ & finite engineering choice\\
8 bytes & 64 & $18{,}446{,}744{,}073{,}709{,}551{,}616$ & $2^{32}$ & current prototype\\
\bottomrule
\end{tabular}
\caption{Fixed fingerprint widths.  These scales do not imply an all-input
collision bound.}
\label{tab:digits}
\end{table}

For example, with $Q=100{,}000$ pairwise distinct states, the expected
numbers of colliding pairs are approximately $76{,}293$ at four hex digits,
$298$ at six, and $1.16$ at eight.  Those differences can be substantial in
a finite implementation.  All three widths remain $O(1)$ as $M\to\infty$,
so none changes a polynomial exponent or defeats the fixed-map adversary.
If the output alphabet is decimal, the birthday scales for four, six and
eight characters are instead roughly $10^2,10^3,10^4$ states.

\section{Adaptive exact fingerprint construction}
The strongest safe engineering variant is a verified multilevel cache.
Start with a short prefix, for example sixteen bits.  When a bucket exceeds a
declared cap, extend only that bucket by another block of digest bits.  Repeat
through twenty-four, thirty-two and sixty-four bits, and finally compare the
exact normalized residual stream before every reuse.  An exact descriptor or
source-replay fallback handles the final level.

\begin{theorem}[ADAPTIVE-EXACT1]
The verified multilevel cache returns exactly the answer returned by the same
solver with caching disabled, for every choice of prefix lengths, bucket caps
and hash function.
\end{theorem}
\begin{proof}
A cache answer is accepted only after the final exact congruence test.
Short prefixes and extensions merely choose the order and number of
candidates.  A miss invokes the uncached exact computation.  Induction over
the solver call graph therefore preserves every returned Boolean value.
\end{proof}

This construction can reduce average tag storage and average replay on
benign inputs.  It cannot improve the deterministic worst case by itself.
For a fixed sequence $H_1,\ldots,H_t$ of $h$-bit hashes, the concatenated
signature has $th$ bits, and among $R$ semantic classes one signature still
contains at least
\[
 \left\lceil R/2^{th}\right\rceil
\]
classes.  Repeated hashing accumulates the information of one longer
fingerprint.  Fresh random rehashing gives a Las Vegas dictionary only when
every reuse remains verified; it gives expected rather than deterministic
resource bounds.  Deterministic seed enumeration must charge failed seeds and
collision certification.

\subsection{CACHE-FINGERPRINT1 bounded test}
The companion program \texttt{cache\_fingerprint\_experiment.py} implements
fixed two-, four-, six- and eight-hexadecimal-digit caches and adaptive
$2\mathord\to8$ and $4\mathord\to8$ variants.  Each prospective hit still
receives a full normalized-residual-stream comparison.  The old residual is
regenerated whenever an adaptive entry needs extension bits, so extension
does not hide source work.  Each Python process has a 25-second alarm.

The final batch used unsatisfiable PHP$(7,6)$ and eight satisfiable
interleaved isomorphic dense components.  It completed twelve answer checks
in 2.90 seconds with no error.  All configurations made the same 892
recursive calls and obtained the same 121 verified cache hits.

\begin{center}
\resizebox{\textwidth}{!}{%
\begin{tabular}{lrrrrrr}
\toprule
scheme & max cache bits & bucket peak & equality replays & rejected collisions & literal reads & peak workspace\\
\midrule
fixed 2 hex & 49,244 & 10 & 1,278 & 1,157 & 887,763 & 51,721\\
fixed 4 hex & 55,396 & 2 & 122 & 1 & 833,920 & 57,577\\
fixed 6 hex & 61,548 & 1 & 121 & 0 & 833,889 & 63,433\\
fixed 8 hex & 67,700 & 1 & 121 & 0 & 833,889 & 69,289\\
adaptive $2\mathord\to8$ & 55,116 & $10^*$ & 312 & 191 & 989,748 & 57,273\\
adaptive $4\mathord\to8$ & 55,412 & $2^*$ & 122 & 1 & 834,328 & 57,593\\
\bottomrule
\end{tabular}}
\end{center}
The starred adaptive values are occupancies before extension; the final
stored-prefix bucket peak was one.  On this finite batch, four hex digits were
the best measured fixed tradeoff.  Moving from four to six removed one false
collision and saved only 31 literal reads while adding 6,152 cache bits.
Moving from six to eight added another 6,152 bits without reducing work.
Starting adaptively at two digits reduced rejected comparisons relative to
fixed two digits, but 712 extension rehashes increased literal reads by 11.5
percent.  These measurements select an engineering setting for these two
workloads; they do not establish an asymptotic optimum or a general SAT
bound.  Complete raw data are in
\texttt{cache\_fingerprint\_results.json}.

\subsection{Collision-focused hostile controls}
A separate finite generator searched for short-digest collisions among
syntactically distinct residuals and then replayed each saved witness.  At
sixteen bits and 100,000 residuals it found 29,520 colliding buckets, 18,688
of which mixed SAT answers, with maximum occupancy nine.  At twenty-four bits
and 250,000 residuals it found 1,796 colliding buckets, 846 cross-answer,
with maximum occupancy three.  At thirty-two bits and 500,000 residuals it
found 36 colliding buckets, 18 cross-answer, with maximum occupancy two.

A cache that merged by the digest alone returned respectively 1,288, 404 and
72 wrong answers on collision-focused four-pass workloads.  Full-key and
verified-hash caches returned no wrong answer.  The verified variants paid
10,426, 1,849 and 1,128 exact equality comparisons and, in these runs, stored
slightly more than the full-key cache.  The control also found equal-current-
answer residuals that became unequal after applying the same later
restriction.  This is a concrete instance of the congruence requirement:
the current SAT bit is not a safe recursive merge key.  The finite counts are
\textbf{COMPUTATIONAL}; the pigeonhole and congruence lemmas supply the
general statements.

\section{Short summaries cannot select a universal restriction cover}
There is a direct guard against using the digest itself to select the small
family of WRAPPER children.

\begin{lemma}[SUMMARY-COVER1]
For every $x\in\{0,1\}^n$, let $F_x$ be the unit-clause CNF with unique
satisfying assignment $x$.  Suppose a constructor sees only an $h$-bit
summary $H(F_x)$ and, for each summary, emits at most $q$ partial assignments,
each leaving at most $\ell$ variables free.  If these restrictions cover the
unique satisfying assignment for every $F_x$, then
\[
 \boxed{h+\log_2q+\ell\ge n.}
\tag{10}
\]
\end{lemma}
\begin{proof}
Some summary bucket contains at least $2^{n-h}$ strings $x$.  The emitted
family is identical for every formula in this bucket.  One restriction with
$\ell$ free variables covers at most $2^\ell$ assignments, so all $q$ cover
at most $q2^\ell$.  Exact coverage requires
$q2^\ell\ge2^{n-h}$, which is (10).
\end{proof}

For the ordinary encoding of these unit formulas, $M=\Theta(n\log n)$.
Taking $h=O(\log M)$, $q=M^{1/3+o(1)}$ and
$\ell\le M^{2/3+o(1)}$ makes the left side $o(n)$, contradicting (10).
A constructor that continues to read and exploit the full source escapes the
lemma.  Its source-dependent acquisition work is then exactly what the
recursive interface must charge.  The lemma rules out summary-only covers;
it does not rule out a stronger source-dependent SAT reduction.

\section{Recovered project results}
The past-work search found four results that constrain the present design.
\begin{description}
\item[] \textbf{PRIME-LINEAGE-CNF1, PROVED/EXACT.} Distinct prime or Boolean
labels can encode syntactic provenance inside a CNF with polynomial bit
length.  They do not force an arbitrary SAT decider to retain the lineage.
\item[] \textbf{Shared BAD-SEED construction, PROVED/EXACT.} Exponentially
many routes can be represented by a polynomial shared circuit and replayed
from a short seed.  This removes route enumeration as an automatic lower
bound, while the final SAT bit can still bypass history retention.
\item[] \textbf{WORKSAT-0011, COMPUTATIONAL.} A rolling $O(\log N)$ stream
fingerprint greatly reduced measured storage relative to a full transcript,
but only fingerprint-weak sameness was obtained; injective strong sameness
and the unrestricted zero-space floor remained unavailable.
\item[] \textbf{WORKSAT-0017/0019, FALSIFIED as general POWER.} A truncated
sixteen-bit lossy hash failed strong sameness.  Combining a digest with an
adaptive pointer either answered equality only, re-solved on a miss, or
rematerialized linear content for unrestricted access.  This is the earlier
version of the present exactness-versus-access boundary.
\end{description}

The continuous random-bitstream proposal was likewise retained only as
source-binding and hostile-test infrastructure.  Without exact verification
it introduces error; with verification it supplies randomized expected
bounds; a favorable fixed stream is advice; and a single globally fixed
uncomputable stream is an oracle-like object.  A computable fixed stream is a
public deterministic schedule and admits hostile inputs.

\section{Literature-guided positive routes}
The closest positive construction found in the current literature uses the
same safe division of labor as this manuscript: a cheap fingerprint proposes
a merge and a source-backed test certifies it.  Biere, Lagniez and Lonca first
hash a decision-DNNF node by its variable and exact model count, then verify a
candidate semantic merge by two provenance-based entailment checks
\cite{bierelagniezlonca2026}.  The verifier is polynomial on the already
compiled representation.  The initial compilation of an arbitrary CNF can
be exponential, the model counts have growing bit length, and hash buckets
can still contain many candidates.  This is therefore an exact realization
of the desired merge interface after acquisition, rather than a proof of
cheap acquisition on general SAT.

Canonical decision diagrams give a second restricted route.  For a fixed
variable order, reduced ordered BDDs are canonical and exact hash-consing
compares the full tuple consisting of a variable and its two child identifiers
\cite{bryant1986}.  Tree Decision Diagrams generalize this pattern.  Capelli
et al. prove polynomial minimization for a fixed bounded-width TDD and
$2^{O(k)}$ compilation bounds for several width-$k$ source classes
\cite{capelli2026}.  They also exhibit functions requiring exponential TDD
size.  A useful WRAPPER target is consequently a source-checkable theorem
that every generated residual remains in one of these bounded-width regimes.
No such theorem is known for arbitrary SAT.

Component caching in DPLL supplies the analogous algorithmic lesson.
Exact residual components can be reused with favorable bounds controlled by
branchwidth, while general traces and compiled decision diagrams can remain
exponential \cite{bacchusdalmaopitassi2014}.  Hash length is an implementation
parameter after the structural parameter is bounded.  It does not establish
that bound.

Adaptive approximate-membership structures append fingerprint bits when
collisions are discovered, but their theorems are randomized, permit false
positives under a specified adversary model, and retain access to the full
keys \cite{benderetal2018}.  They support the adaptive cache experiment as an
engineering design, not as an exact semantic quotient.

\section{Relation to the proof spines}
In \texttt{main5.tex}, changing $h$ can improve the implementation of W3 and
W6 after a small exact class set has been proved.  HASH-INDEX-NO-SHRINK1 says
that it does not establish W2, W4 or W5.  In \texttt{main6.tex}, the same
point says that a hash does not bound the forward-tree node count $Q$ or leaf
length $M(N)$ unless a source-checkable semantic theorem already connects
the buckets to those objects.  In \texttt{main4.tex}, finite constant-factor
savings do not create a new achievable exponent pair.

To obtain at most $M^{1/3+o(1)}$ children by declaring one child per hash
bucket, one would naturally choose
$h\le(1/3+o(1))\log_2M$.  Exact verification then leaves the child count at
$R$, which the hash does not bound.  Merging all equal hashes instead requires
proving (1) and constructing the semantic quotient within the target work and
space bounds.  This is COVER-SHRINK1 itself.

The strongest justified use is therefore:
\[
\begin{gathered}
\text{source-dependent theorem proves few exact canonical classes}\\
\Longrightarrow\quad
\text{adaptive fingerprint indexes those classes cheaply}\\
\Longrightarrow\quad
\text{exact replay preserves correctness.}
\end{gathered}
\]
The first arrow remains open on general SAT.  Four, six and eight digits are
useful experimental settings for the second arrow, but none proves the first.
P versus NP remains open.

\begin{thebibliography}{9}
\bibitem{carterwegman1979}
J. Lawrence Carter and Mark N. Wegman,
\emph{Universal Classes of Hash Functions},
Journal of Computer and System Sciences 18(2):143--154, 1979.
\url{https://doi.org/10.1016/0022-0000(79)90044-8}

\bibitem{karprabin1987}
Richard M. Karp and Michael O. Rabin,
\emph{Efficient Randomized Pattern-Matching Algorithms},
IBM Journal of Research and Development 31(2):249--260, 1987.
\url{https://doi.org/10.1147/rd.312.0249}

\bibitem{fks1984}
Michael L. Fredman, J\'anos Koml\'os and Endre Szemer\'edi,
\emph{Storing a Sparse Table with $O(1)$ Worst Case Access Time},
Journal of the ACM 31(3):538--544, 1984.
\url{https://doi.org/10.1145/828.1884}

\bibitem{bierelagniezlonca2026}
Armin Biere, Jean-Marie Lagniez and Emmanuel Lonca,
\emph{Efficient Minimization of Decision-DNNF Circuits via Semantic Hashing
and Provenance Tracking}, IJCAI 2026, pages 3774--3782.
\url{https://www.ijcai.org/proceedings/2026/420}

\bibitem{bryant1986}
Randal E. Bryant,
\emph{Graph-Based Algorithms for Boolean Function Manipulation},
IEEE Transactions on Computers C-35(8):677--691, 1986.
\url{https://doi.org/10.1109/TC.1986.1676819}

\bibitem{capelli2026}
Florent Capelli, YooJung Choi, Stefan Mengel, Mart\'in Mu\~noz and
Guy Van den Broeck,
\emph{A Canonical Generalization of OBDD}, SAT 2026.
\url{https://doi.org/10.4230/LIPIcs.SAT.2026.10}

\bibitem{bacchusdalmaopitassi2014}
Fahiem Bacchus, Shannon Dalmao and Toniann Pitassi,
\emph{Solving \#SAT and Bayesian Inference with Backtracking Search},
Journal of Artificial Intelligence Research 34, 2009; manuscript revision
available as arXiv:1401.3458.
\url{https://arxiv.org/abs/1401.3458}

\bibitem{benderetal2018}
Michael A. Bender et al.,
\emph{Bloom Filters, Adaptivity, and the Dictionary Problem},
FOCS 2018.
\url{https://arxiv.org/abs/1711.01616}
\end{thebibliography}
\end{document}
